\documentclass[11pt]{amsart}
\usepackage{amsmath,amssymb}
\usepackage{graphicx}
\usepackage[hidelinks]{hyperref}

\newcommand{\R}{\mathbb{R}}
\newcommand{\C}{\mathbb{C}}
\newcommand{\Om}{\Omega}
\newcommand{\Wzero}{W^{1,2}_0(\Omega)}
\newcommand{\tform}{\mathfrak{t}}

\title[Constant Dirichlet Angle Characterization]{A constant-coefficient Dirichlet case of the angle-coincidence question}
\author{Agent \texttt{agent\_super\_00}}
\date{2026-07-01}

\begin{document}
\maketitle

\section*{Source and Status}

This packet concerns arXiv:2510.23450, Meinlschmidt--Rehberg,
\emph{Sharp angle estimates for second order divergence operators}~\cite{MR}.
In the final comments, the source asks for useful information on when the
inequality
\[
  \omega(\tform_V)\leq \omega(\mu)
\]
is sharp. The source gives a two-dimensional constant skew-symmetric
Dirichlet example where the inequality is strict, and then says that a precise
useful characterization seems hard.

\begin{center}
\includegraphics[width=0.95\textwidth]{figures/open_problem_crop.png}
\end{center}

The result below supplies a small but exact characterization in the constant
real full-Dirichlet class.

\section*{Definitions}

Let $\Omega\subseteq\R^d$ be a nonempty open set and set
$V=\Wzero$. For a constant real matrix $M\in\R^{d\times d}$ whose symmetric
part is uniformly positive, consider the Dirichlet form
\[
  \tform(u,v)=\int_\Omega (M\nabla u,\nabla v)_{\C^d}\,dx,
  \qquad u,v\in \Wzero.
\]
Write
\[
  M=S+K,\qquad S=\frac{M+M^T}{2},\qquad K=\frac{M-M^T}{2}.
\]
Thus $S=S^T$ and $K^T=-K$. Assume $S\geq mI$ for some $m>0$.

\section*{Proof Intuition}

With full Dirichlet boundary conditions, a constant skew-symmetric coefficient
matrix contributes only a divergence-free second-order expression. Integration
by parts kills it in the weak form. The form angle is therefore controlled
only by the symmetric positive part $S$, so it is zero. The coefficient angle,
however, is computed from the numerical range of the matrix itself; a nonzero
real skew-symmetric part produces a nonzero imaginary numerical direction.
This separates the two angles exactly.

\section*{Theorem}

Let $\Omega\subseteq\R^d$ be a nonempty open set and let
$M=S+K\in\R^{d\times d}$ be constant with $S=S^T\geq mI$ and $K^T=-K$. For the
full Dirichlet form on $V=W^{1,2}_0(\Omega)$,
\[
  \omega(\tform_V)=0.
\]
Moreover, the coefficient angle satisfies $\omega(M)=0$ if and only if
$K=0$. Consequently, in this class,
\[
  \omega(\tform_V)=\omega(M)
  \quad\Longleftrightarrow\quad
  M \text{ is symmetric}.
\]

\begin{proof}
First suppose $u,v\in C_c^\infty(\Omega)$. Since $K$ is constant and
skew-symmetric,
\[
  \int_\Omega (K\nabla u,\nabla v)_{\C^d}\,dx
  =
  \sum_{j,k=1}^d K_{jk}\int_\Omega \partial_k u\,\overline{\partial_j v}\,dx.
\]
Integrating by parts in the $j$ variable gives
\[
  -\sum_{j,k=1}^d K_{jk}\int_\Omega
  \partial_j\partial_k u\,\overline{v}\,dx=0,
\]
because the matrix $(\partial_j\partial_k u)_{j,k}$ is symmetric while
$K$ is skew-symmetric. By density of $C_c^\infty(\Omega)$ in $W^{1,2}_0(\Omega)$,
the identity
\[
  \int_\Omega (K\nabla u,\nabla v)_{\C^d}\,dx=0
\]
holds for all $u,v\in W^{1,2}_0(\Omega)$. Hence
\[
  \tform(u,u)=\int_\Omega (S\nabla u,\nabla u)_{\C^d}\,dx\in [0,\infty)
\]
for every $u\in W^{1,2}_0(\Omega)$, and therefore the numerical range of the
form is contained in the positive real axis. Thus $\omega(\tform_V)=0$.

It remains to identify when the coefficient angle is zero. If $K=0$, then
$M=S$ is selfadjoint and positive, so $N(M)\subset [0,\infty)$ and
$\omega(M)=0$. Conversely, assume $K\neq 0$. The operator
\[
  \operatorname{ImOp}(M)=\frac{M-M^T}{2i}=-iK
\]
is a nonzero selfadjoint operator on $\C^d$, hence has a nonzero real
eigenvalue $\lambda$ with unit eigenvector $\xi$. Then
\[
  \operatorname{Im}(M\xi,\xi)_{\C^d}
  =(\operatorname{ImOp}(M)\xi,\xi)_{\C^d}
  =\lambda\neq 0,
\]
whereas
\[
  \operatorname{Re}(M\xi,\xi)_{\C^d}
  =(S\xi,\xi)_{\C^d}>0.
\]
Thus $N(M)$ is not contained in the positive real axis, and $\omega(M)>0$.
This proves the equivalence.
\end{proof}

\section*{Verification Notes and Scope}

The proof is purely analytic. It uses only density of $C_c^\infty(\Omega)$ in
$W^{1,2}_0(\Omega)$ and the elementary spectral fact that a nonzero
selfadjoint finite-dimensional operator has a nonzero numerical direction.

The result is intentionally narrow. It gives a complete characterization in
the constant real full-Dirichlet class, extending the source paper's displayed
two-dimensional skew example, but it does not handle variable coefficients,
mixed or Neumann boundary conditions, or the complex-coefficient
$p$-numerical-range problem raised in the next item of the source comments.

\section*{Human Review Recommendation}

Review as a scoped partial result toward the angle-coincidence question. The
main point to check is the integration-by-parts cancellation of the constant
skew-symmetric part in the full Dirichlet form.

\begin{thebibliography}{9}
\bibitem{MR}
H. Meinlschmidt and J. Rehberg,
\emph{Sharp angle estimates for second order divergence operators},
arXiv:2510.23450.
\end{thebibliography}

\end{document}
